
\section{Discussion}\label{sec:disc}

\subsection{What the design achieved}
Moving the patient's decisions into code made its behaviour something an examiner can watch, a reviewer can test and a researcher can argue
about. Safety rules became automatic gates, and an independent reviewer could try to break them, and did; all ten defects it found were fixed
and turned into tests before the results were measured. The ablation shows where the safety result comes from: without the decision rubric the
patient failed four jailbreak attempts and fell back to the old, narrow pattern of behaviour (Table~\ref{p1:tab:ablation}). The patient's behaviour
now changes with the kind of doctor it faces, while each single conversation stays as focused as before. Because the patient is data, a new
station is a new file rather than new code, and because the body reads the same decision as the voice, posture, gaze and face agree with what
the patient says. All of this runs on one laptop, so no student's voice leaves the device.

\subsection{What fell short}
\textbf{Accuracy did not rise.} The composite measures accuracy against the case, not realism, and the previous Twin was already accurate.
Whether the added fit to the situation feels more real is a question for clinicians; a blind A/B page with ten pairs of replies is ready but not
yet rated.

\textbf{Recall of reworded risk questions} reached 72.7\% on the held-out half, against a 95\% target. The same method scored 90.7\% on the tuning
half, so a tuning half that is also used to write examples gives a badly over-optimistic number. Tuning only to find more risk questions made the
patient talk about suicide when asked about sleep, which is why recall must be tuned with a limit on wrong matches.

\textbf{An angry patient calmed down too soon.} The upset style was in character on half of the turns after turn 10, because the softening rules
let a consistently kind doctor lower resistance. A minimum hold time for upset patients is the obvious fix.

\textbf{Speed.} A median of 1.8\,s is close to the thoughtful pause of a low-mood patient, but misses a 1.2\,s target for brisker cases. Holding
a whole reply back on a risk turn, so that a denial can be replaced, and the longer prompt both cost time; part of the increase is not yet
traced.

\textbf{Rules that read the doctor.} Word patterns miss some judgemental sentences, such as ``Are you sure you're not just making this up?'', and
the rule against a flat ``No'' on locked topics also fired for an accusation about a non-risk topic, where a plain ``No'' would be natural
(Table~\ref{p1:tab:example}).

\subsection{Lessons for builders}
\begin{enumerate}
\item Keep everything a clinician would audit outside the language model, and give the model a closed list to choose from.
\item Check every sentence before it is spoken, and hold back a whole reply when a single sentence could break a safety rule.
\item Write recall test questions \emph{blind}, from what the memory means, never from the matching rules; keep a test half that is never
touched, and report it rather than the tuning half.
\item Always report wrong risk matches next to risk recall, and use a stricter ``clearly asked'' level for anything with consequences, such as
forcing a disclosure.
\item Let a learning loop change only how the patient speaks, never the case, and gate every learned version on the held-out safety checks; in
our first round the loop was most useful as a test that found bugs.
\item Test with a person early: one playtest found four problems that every scripted test had missed.
\end{enumerate}

\subsection{Limitations and threats to validity}
\begin{itemize}
\item Only four draft stations; psychosis, mania and anxiety are not covered. All clinical content, trust levels and debrief cut-offs are drafts
that psychiatrists must still review.
\item The held-out conversations and the examination tests were written by the developer, and the recall benchmark by a single AI writer, whose
phrasing may be more regular than a nervous student's. The tuning and test halves come from the same writer, so the test half measures new
wording, not a new writer.
\item The judge is a language model and its flags were checked by the developer, not by faculty. LLM judges are known to be biased
\cite{zheng2023judging}.
\item \emph{Construct:} the composite measures accuracy, not realism, and the variety measures count labels, not how a reply feels.
\emph{Statistical:} 36 and 6 runs per version detect large differences only; about 75 risk pairs per benchmark half make each recall figure
uncertain by several points.
\item \emph{External:} two cases in the held-out runs, scripted doctors, one variety of English, one high-end laptop. The student's own body
language is not assessed.
\item The Twin is inspired by brain science but is not a model of the brain.
\end{itemize}

\subsection{Ethics}
No human participants or real patient data were involved. A student's words never become patient facts, and student sessions are used for
learning only with consent. Suicide-related content comes with a content note and a debrief with teaching points. Patient voices are designed
from written descriptions, not copied from real people. Release builds refuse unreviewed clinical content, and demonstrations that use
pre-scripted replies must say so.

\subsection{Next steps}
Table~\ref{tab:next} lists the next steps, in order.

\begin{table}[t]
\caption{Next steps}
\label{tab:next}
\centering
\scriptsize
\begin{tabularx}{\columnwidth}{@{}YY@{}}
\toprule
Step & Why\\
\midrule
Faculty review of all four cases and the flagged replies & Clinical truth and judge accuracy\\
Faculty blind A/B ratings & Answer RQ3: does it feel real?\\
Ethics approval and a student pilot (n $\approx$ 20--30) & Answer RQ5: does it help?\\
Hand-checked labels and a fine-tuned recall model & Reach 95\% recall on risk questions\\
Learned move rules and a hold time for upset patients & Fewer missed judgements; anger that lasts\\
A psychosis or anxiety case & Broader coverage\\
Lip sync and skin, eye and hair shading & A more lifelike face up close\\
\bottomrule
\end{tabularx}
\end{table}

\section{Conclusion}
A spoken, embodied psychiatric standardised patient can run privately on one laptop if the patient is data, code makes the decisions, a local
language model supplies the words, and the body reads the same decision as the voice. With this split, safety becomes a set of checks that anyone
can run and try to break, and the remaining gaps, such as recall of risk questions asked in new words and keeping an angry patient angry, become
measurable targets instead of impressions. Whether the patient is real enough, and useful enough, to teach with is the next question, and it
needs psychiatrists and students to answer it.

\appendices
\section{What a seed changes}
A seed picks one option for each everyday ``slot'' in the case and one conversation style allowed by the case. For the depression case the slots
include a hobby (for example Carnatic singing or badminton), a younger brother's name and a home town; every option is written by hand, so
nothing clinical can change. The answer key, the memories' trust levels, the rating-scale scores and the risk items are the same for every seed.
Assessment mode uses the authored style for every student so that marks are comparable; practice mode allows any style.

\section{Instructions given to the benchmark writer}
The writer received, for each case: the patient's name, age and situation; for each memory its id, its kind (complaint, symptom, history, risk,
social, belief) and a one-line description of what the patient knows; and these rules: write lines a real doctor might say to this patient;
at least 4 per memory and 12 per risk memory; mix textbook, gentle and indirect, clumsy and leading, colloquial Indian English, and euphemistic
phrasings; most lines should avoid the memory's obvious key words; each line must make sense on its own; label each line with every memory it
asks about, and nothing more; include greetings, small talk and empathy-only lines that ask about nothing; mark hard lines. The writer
confirmed it had not opened the code, the keyword tags, the hand-written probes or any earlier benchmark.

\section{Hand-set heads: example questions}
Each head is the normalised average of its example questions' vectors. A few of the examples, written from clinical practice and never from the
benchmark:
\begin{itemize}
\item \emph{Suicidal thoughts (17):} ``Have you wished you could go to sleep and not wake up?'', ``Do you sometimes feel everyone would be better
off without you?'', ``When it gets very heavy, do thoughts of death come to you?''
\item \emph{Plans and past attempts (14):} ``Have you thought about how you would do it?'', ``Have you done anything to prepare, like writing a
note?'', ``Can you keep yourself safe if you go home today?''
\item \emph{Means (12):} ``Have you been saving up medicines?'', ``Can you get onto a roof or somewhere high easily?'', ``Is your home safe for
you right now?''
\item \emph{Harm to others:} ``Has anyone at home been afraid of you?'', ``Have arguments ever got physical?''
\end{itemize}

\section{System requirements and reproducing the results}
\textbf{Hardware and models.} Apple M5 Pro, 24\,GB memory. Patient: Gemma 4 E4B, 4-bit (Q4\_K\_M), llama.cpp server with Metal, flash
attention, all layers on the graphics chip, two parallel slots. Judge: Qwen3.5-9B on a second server. Recall: bge-small-en-v1.5, ONNX, on the
processor. The voice model alone peaked at 8.1\,GB in the bake-off, so a machine with less memory may not fit the expressive voice alongside the
language model.

\textbf{Held-out runs.} For each version $v\in\{\text{v0},\text{v2}\}$, case and seed, run
\texttt{tools/learning/loop.py heldout $v$ --case $c$ --seeds 11,23,42} with the judge URL; \texttt{papers/analysis/analyses.py} computes the
pooled and per-conversation numbers, intervals and decision counts, and \texttt{papers/figures/make\_figures.py} draws the figures.

\textbf{Recall.} \texttt{tools/recall\_bench.py --split test} prints recall and wrong-match rates for the shipped settings; \texttt{--lexical}
gives the keyword baseline and \texttt{--hand} adds the hand-set heads. \texttt{papers/analysis/analyses.py --recall} and
\texttt{recall\_extra.py} compute the by-case, difficulty, probe-source, rank, interval and timing results.

\textbf{Speed.} \texttt{tools/latency/ws\_bench.py} plays synthesised Indian-English doctor lines into the real session connection as 32\,ms
frames and never uses the microphone; results come from the service's own timing log.

The held-out scripts, case files, brain versions, benchmark and its split, and every judged reply are kept with the code, so each number in this
article can be traced to the data behind it.

\section*{Acknowledgment}
\textbf{Use of AI tools.} The text of this article, its figures, tables and analysis scripts, and the \LaTeX{} source were drafted with an AI
assistant (Claude, Anthropic, used through Claude Code), working from the project's code, experiment logs and result files under the author's
direction. The same assistant helped write the software described here, ran the analyses, and checked the citations against their sources.
Within the study itself, AI models were also used as research tools, as described in Section~\ref{sec:eval}: a separate language model
(Qwen3.5-9B) marked the held-out replies, an AI agent wrote the blind recall benchmark, and another AI agent carried out the independent safety
review. The author takes full responsibility for the content of this article.

\bibliographystyle{IEEEtran}
\bibliography{../shared/refs}

\end{document}
